% !PRINTED 71
\weberpage{71}
%<<CONT
gleich einem neuen $\vartheta$ setzt, so ergibt sich f\"ur dieses die partielle
Differentialgleichung

\weq{4}{\frac{\partial^2 \vartheta(v,\omega)}{\partial v^2} - 4\pi i\,\frac{\partial \vartheta(v,\omega)}{\partial \omega} = 0,}

\gesperrt{und durch (1) und (4) ist jetzt die Funktion $\vartheta$ bis auf
einen von $v$ und $\omega$ unabh\"angigen, also nur noch von $g_1, g_2$
abh\"angigen Faktor definiert.}

\begin{center}
\textbf{\S\,21.\quad\gesperrt{Die Theta-Funktionen verschiedener Charakteristiken.}}

\textbf{\gesperrt{Hauptcharakteristiken.}}
\end{center}

Wir wollen jetzt, ehe wir an die Bestimmung des noch
\"ubrigen von $v$ und $\omega$ unabh\"angigen Faktors in den $\vartheta$-Funktionen
gehen, die verschiedenen Charakteristiken aufeinander zur\"uck-
f\"uhren, was nach \S\,17 immer m\"oglich ist. Bezeichnet man mit
$\vartheta(v)$ die zur Charakteristik $(0, 0)$ geh\"orige $\vartheta$-Funktion, so ist
\[
\Phi(v,\omega) = e^{-\pi i g_1 v}\, \vartheta\!\left(v - \frac{g_1\omega - g_2}{2}\right)
\]
eine den Bedingungen (1) \S\,20 gen\"ugende Funktion. Es ist aber
nun mit R\"ucksicht auf die Differentialgleichung (4):
\[
\frac{\partial^2 \Phi}{\partial v^2} - 4\pi i\,\frac{\partial \Phi}{\partial \omega} = -\pi^2 g_1^2\,\Phi,
\]
und wenn wir also

\weq{1}{\vartheta_{g_1,\, g_2}(v) = e^{\frac{\pi i \omega g_1^2}{4} - \pi i g_1 v + \frac{g_1 g_2 \pi i}{2}}\ \vartheta\!\left(v - \frac{g_1\omega - g_2}{2}\right)}

\gesperrt{setzen, so ist die Differentialgleichung (4) \S\,20 durch alle
diese Funktionen befriedigt. Hierdurch also sind die Funktionen
 $\vartheta_{g_1,\, g_2}$ bis auf einen allen gemeinschaftlichen numerischen
Faktor bestimmt.}

Aus der Formel (1) l\"a\ss t sich eine allgemeinere ableiten,
indem man $v$ durch $v - \dfrac{g_1'\omega - g_2'}{2}$ ersetzt. Man erh\"alt so

\weq{2}{\vartheta_{g_1,\, g_2}\!\left(v - \frac{g_1'\omega - g_2'}{2}\right)
= e^{-\frac{\pi i \omega}{4}g_1'^2 + \pi i g_1' v - \pi i g_1 g_2' - \frac{\pi i}{2}g_1'(g_2+g_2')}\ \vartheta_{g_1+g_1',\, g_2+g_2'}(v).}
